\clearpage
\chapter*{LAMPIRAN B. \emph{QUERY MEMORY ITEMS}}
\addcontentsline{loa}{appendix}{Lampiran B. Kueri \emph{memory items}}

Lampiran ini memuat kueri untuk mengekstrak \emph{memory item} dari konteks
percakapan. Placeholder \texttt{previous\_items}, \texttt{summary},
\texttt{messages}, dan \texttt{max\_items} diisi oleh sistem saat \emph{runtime}.

\begin{quote}\small\ttfamily\raggedright
Extract important session memory items from the conversation context below.\par
\par
PREVIOUS MEMORY ITEMS:\par
\{prev\_text\}\par
\par
ROLLING SUMMARY:\par
\{summary\}\par
\par
NEW MESSAGES:\par
\{messages\}\par
\par
Return at most \{max\_items\} important memory items.\par
Each item must be short and specific.\par
Keep still-relevant existing items.\par
Remove duplicates.\par
Remove completed or contradicted items.\par
Sort by importance.\par
Use bullet points starting with "- ".\par
Focus on: hard constraints, explicit decisions, current goal,\par
open tasks, corrections, and technical requirements.
\end{quote}

Kueri ini menggabungkan \emph{memory item} sebelumnya, \emph{rolling summary}, dan pesan
terbaru. Hasilnya dibatasi paling banyak \texttt{max\_items} \emph{item}, setiap \emph{item}
harus singkat, spesifik, tidak duplikat, dan tetap relevan. Item yang sudah
selesai atau bertentangan dengan informasi baru harus dikeluarkan.
